% Part: first-order-logic
% Chapter: introduction
% Section: sentences

\documentclass[../../../include/open-logic-section]{subfiles}

\begin{document}

\olfileid{fol}{int}{snt}

\olsection{\usetoken{P}{sentence}}

సరే, ఇప్పుడు \tetoken{నిర్మాణాలు}{!!{structure}s}, \tetoken{సూత్రాలు}{!!{formula}s} కోసం సంతృప్తి
(``దానిలో సత్యం'') యొక్క ఒక (స్థూల) నిర్వచనం ఉంది. కానీ దానికి
\tetoken{ఒక చరం}{!!a{variable}} నిర్దేశం అనే అదనపు అంశం అవసరమైంది; మనం కోరుకున్నది
\tetoken{వాక్యాలు}{!!{sentence}s} నిర్వచనం. నిర్దేశాలను ఎలా తొలగిస్తాం, అసలు
\tetoken{వాక్యాలు}{!!{sentence}s} అంటే ఏమిటి?

\tetoken{చరాలకు}{!!{variable}s}, పరిమాణీకరణికాలకు మధ్య సంబంధం గురించి ఔపచారిక
తర్కానికి మీ మొదటి పరిచయంలో జరిగిన చర్చ బహుశా మీకు గుర్తుండవచ్చు.
ఒక పరిమాణీకరణిక తరువాత ఎప్పుడూ ఒక \tetoken{చరం}{!!a{variable}} వస్తుంది; ఆ
పరిమాణీకరణిక వర్తించే \tetoken{వాక్యం}{!!{sentence}} భాగంలో (దాని ``పరిధి''లో), ఆ
\tetoken{చరం}{!!{variable}} ఆ పరిమాణీకరణిక చేత ``బద్ధం'' అయిందని భావిస్తాం.
\tetoken{సూత్రాలలో}{!!{formula}s} ప్రతి \tetoken{చరానికి}{!!{variable}} సరిపోలే పరిమాణీకరణిక ఉండాలని
కోరలేదు; సరిపోలే పరిమాణీకరణికలు లేని \tetoken{చరాలు}{!!{variable}s} ``స్వేచ్ఛా'' లేదా
``బంధితం కాని'' చరరాశులు. స్వేచ్ఛా \tetoken{చరాలు}{!!{variable}s} ఏవీ లేని \tetoken{సూత్రాలు}{!!{formula}s}
అన్నింటినీ \tetoken{వాక్యాలుగా}{!!{sentence}s} తీసుకుంటాం.

మళ్లీ, \tetoken{ఒక సూత్రం}{!!a{formula}}~$!A$లో \tetoken{ఒక చరం}{!!a{variable}} యొక్క ఒక సంభవం ఎప్పుడు
బద్ధం, దాన్ని ఏ పరిమాణీకరణిక బంధిస్తుంది, అది ఎప్పుడు స్వేచ్ఛా సంభవం
అనే అంతర్బుద్ధి గ్రహించడం కష్టం కాదు. దీన్ని పరీక్షించడానికి,
బహుశా కుండలీకరణలను లెక్కించే ఒక పద్ధతిని మీరు నేర్చుకొని ఉండవచ్చు.
మనం మాత్రం ఖచ్చితమైన నిర్వచనాన్నే కోరాలి. \tetoken{సూత్రాలను}{!!{formula}s} ఆగమనం
ద్వారా నిర్వచించాం కాబట్టి, \tetoken{ఒక సూత్రం}{!!a{formula}}~$!A$లో \tetoken{ఒక చరం}{!!a{variable}}~$x$
యొక్క స్వేచ్ఛా, బద్ధ సంభవాలను కూడా ఆగమనం ద్వారా నిర్వచించవచ్చు.
ఉదాహరణకు, మన సరళీకృత భాషలో అది ఇలా ఉండవచ్చు:
\begin{enumerate}
  \item $!A$ పరమాణు సూత్రమైతే, దానిలోని $x$ సంభవాలన్నీ స్వేచ్ఛా
  సంభవాలే (అంటే, $\Atom{\Obj P}{x}$లోని $x$ సంభవం స్వేచ్ఛా సంభవం).
  \item $!A$ రూపం $\lnot !B$ అయితే, $\lnot !B$లోని $x$ యొక్క ఒక
  సంభవం స్వేచ్ఛా సంభవం కావడం, $!B$లో దానికి సరిపోలే $x$ సంభవం
  స్వేచ్ఛా సంభవమైనప్పుడూ, అప్పుడు మాత్రమే (అంటే, $!B$లోని చరరాశుల
  స్వేచ్ఛా సంభవాలే $\lnot !B$లోని వాటికి సరిపోలే సంభవాలు).
  \item $!A$ రూపం $(!B \land !C)$ అయితే, $(!B \land !C)$లోని $x$
  యొక్క ఒక సంభవం స్వేచ్ఛా సంభవం కావడం, $!B$లో లేదా $!C$లో దానికి
  సరిపోలే $x$ సంభవం స్వేచ్ఛా సంభవమైనప్పుడూ, అప్పుడు మాత్రమే.
  \item $!A$ రూపం $\lexists[x][!B]$ అయితే, $!A$లోని $x$ సంభవం ఏదీ
  స్వేచ్ఛా సంభవం కాదు; అది $\lexists[y][!B]$ రూపంలో ఉండి, $y$ అనేది
  $x$కు భిన్నమైన \tetoken{ఒక చరం}{!!{variable}} అయితే, $\lexists[y][!B]$లోని $x$
  యొక్క ఒక సంభవం స్వేచ్ఛా సంభవం కావడం, $!B$లో దానికి సరిపోలే $x$
  సంభవం స్వేచ్ఛా సంభవమైనప్పుడూ, అప్పుడు మాత్రమే.
\end{enumerate}

చరరాశుల స్వేచ్ఛా, బద్ధ సంభవాలకు ఖచ్చితమైన నిర్వచనం లభించిన తరువాత,
స్వేచ్ఛా \tetoken{చరం}{!!{variable}s} సంభవాలు లేని ఏ \tetoken{ఒక సూత్రం}{!!{formula}} అయినా
\tetoken{ఒక వాక్యం}{!!a{sentence}} అని సరళంగా చెప్పవచ్చు.

\end{document}
